\documentclass[]{puyang-resume}
\fullname{Pu Yang}
% \jobtitle{Sofware Engineer}

\begin{document}
\resumeheader
{\email{puyang@uvic.ca}}
{\github{puyanguvic}}
{\phone{+1 778-967-2699}}
{\website{https://puyang.me/}}
{\nationality{Canada}}

\begin{section}{Summary}
Senior Machine Learning Researcher focused on trustworthy AI, secure LLM agents, and AI-for-security systems.\\
Lead a team of four Ph.D. researchers on controllable security data synthesis; also research AI-agent permission management and federated-learning security.\\
Research spans LLM safety, agent governance, efficient inference, distributed AI systems, and networked security.
\end{section}

\sectiontable{Technical Skills}{
    \entry{AI/ML}{Python, PyTorch, Transformers, LLM agents, SLM adaptation, diffusion models, federated learning}
    \entry{Systems}{KV-cache optimization, model serving, distributed systems, networking, low-latency inference}
    \entry{Security}{AI firewall, adversarial robustness, model poisoning, backdoors, content safety, policy enforcement}
}

% \begin{section}{Research Experience}
%      \begin{subsection}{Wireless Avionics Intra-Communications}{}{Dec. 2018 -- Dec. 2019}{}
%         \item Sub-6G (4 200 MHz - 4 400 MHz)channel modeling
%     \end{subsection}
% \end{section}

\begin{section}{Work Experience}
    \begin{subsection}{Edison Research Center, Huawei Canada}{Senior Machine Learning Researcher}{Jan. 2026 -- Present}{Canada}
        \item Lead a team of four Ph.D. researchers developing controllable data synthesis methods for AI security applications
        \item Design a GAN-inspired framework that pairs a diffusion language model generator with an encoder-only discriminator to steer generated data toward target styles
        \item Develop permission-management mechanisms for AI agents to govern tool access and resource usage across agentic workflows
        \item Investigate model-poisoning attacks and defense strategies in federated learning to improve robustness against malicious clients
    \end{subsection}

    \begin{subsection}{Edison Research Center, Huawei Canada}{Machine Learning Researcher}{Dec. 2024 -- Dec. 2025}{Canada}
        \item Developed an AI web application firewall using compact BERT-family models to detect malicious and policy-violating inputs at the application boundary
        \item Trained domain-specific tokenizers and small language models for accurate, low-latency security classification
        \item Designed a multimodal phishing-email detection agent that extracts evidence from message content, attachments, and embedded URLs, then orchestrates LLM calls for efficient threat assessment
    \end{subsection}

    \begin{subsection}{6G Network Architecture Design}{Ph.D student}{Dec. 2020 -- Jul. 2023}{Victoria, Canada}
        \item Self-Evolving and Transformative Protocol Architecture for 6G design
        \item Updatable eXpress Transport Protocol design
        \item Delay-Guaranteed Routing Protocol design
    \end{subsection}

    \begin{subsection}{Computing Research Institute of Aviation Industry Corporation of China}{Researcher}{Dec. 2018 -- Dec. 2019}{Xi'an, China}
        \item Sub-6G (4 200 MHz - 4 400 MHz)channel modeling
        \item Wireless communication system design based on Universal Software Radio Peripheral (USRP)
        \item Excellence Research Intern Award Recipient
    \end{subsection}
\end{section}

% \begin{section} {Intern experience}
%     \begin{subsection}{Sichuan Huadi Information Technology Co., Ltd}{Software Development Engineer}{Sep. 2015 -- Oct. 2015}{Chengdu, China}
%     \item Web Development Project Manager for Sichuan Daily
%     \end{subsection}
% \end{section}

\begin{section}{Talks}

\begin{subsection}{Toward the 6G network architecture}{}{July 2022}{Wireless Networking Technology Forum Huawei Canada, Victoria, Canada}
 \item Presented key design principles and challenges for the development of 6G networks, including ultra-low latency, massive connectivity, and enhanced security. Discussed the role of AI-driven network management and potential use cases such as smart cities and autonomous systems.
 \end{subsection}
 
 \begin{subsection}{Multi-path routing protocol to support Quality of Service}{}{Sep, 2023}{IEEE MASS 2023, Toronto, Canada}
 \item Discussed methods to improve network stability in case of congestion and network jamming through the application of multi-path routing protocols, emphasizing their ability to maintain service quality under varying traffic conditions.
 \end{subsection}

 \begin{subsection}{Deadline Driven Routing Protocol}{}{May 2024}{IEEE Conference on Computer Communications, Vancouver, Canada}
 \item Presented the Deadline Driven Routing (DDR) protocol, designed to guarantee service with strict delay constraints in real-time communication systems. Discussed the design, implementation, and evaluation of DDR, focusing on its performance in achieving delay guarantees and improving reliability in delay-sensitive applications.
 \end{subsection}

 \begin{subsection}{Enhancing Cloud Computing with Networking Technique}{}{July 2024}{GenAI and Future Network Forum (Ericsson), Victoria, Canada}
 \item Explored the integration of advanced networking techniques to enhance the efficiency and scalability of cloud computing infrastructures. Demonstrated how novel routing protocols and optimization strategies can improve cloud service delivery, particularly in large-scale distributed environments.
 \end{subsection}
\end{section}




\sectiontable{Professional Service}{
    \entry{Reviewer}{IEEE INFOCOM, IEEE/ACM Transactions on Networking (ToN), Computer Networks, IEEE Internet of Things Journal, and ACM SIGKDD (KDD)}
}

\sectiontable{Awards \& Scholarships}{
    \entry{2024}{Chinese Government Award for Outstanding Self-financed Students Abroad}
    \entry{2020--2023}{UVic Graduate Fellowship; Excellent Research Intern Achievement}
    \entry{2017--2019}{NWPU Graduate Scholarship; ACM-ICPC Second Prize of NWPU (2013)}
}

\begin{section}{Selected Publications}
    \begin{itemize}
        \item \textbf{P. Yang}, T. Chang, L. Cai, “DDR: A Deadline-Driven Routing Protocol for Delay Guaranteed Service” IEEE International Conference on Computer Communications (INFOCOM), Jan 15. 2024.
        \item \textbf{P. Yang}, L. Cai, T. Chang, “DGR: Delay-Guaranteed Routing Protocol” The 20th IEEE International Conference on Mobile Ad-Hoc and Smart Systems (MASS), July. 2023.
        
        \item X. Ren, L. Cai, \textbf{P. Yang}, and J. Ji, “Congestion-aware Delay-guaranteed Scheduling and Routing with Renewal Optimization,” Computer Networks, June 2023.

        \item X. Fan , S. Li , W. Cui , \textbf{P. Yang}, G. Huang, "Mixed-Numerology Channel Division for Wireless Avionics Intracommunications", Wireless Communications and Mobile Computing, April. 2022

    \end{itemize}
\end{section}
\end{document}
